\section{Related Works}
\label{sec:related-work}

The most common evaluation setup for summarization research is to use metrics such as ROUGE~\cite{Lin2004ROUGEAP} or BLEU~\cite{Papineni2002BleuAM} that measure token overlap with a gold reference. These metrics are favored for their interpretability and established use in summarization research, allowing for easy comparison with past work. However, these metrics have a number of well documented issues, such as showing low correlation with human judgment~\cite{Liu2008CorrelationROUGE,fabbri-etal-2021-summeval}. Therefore, these metrics are often combined with other measures, such as BERTScore, which uses a masked language model to measure semantic overlap with a reference summary~\cite{zhang2020bertscoreevaluatingtextgeneration}. Researchers have focused on designing metrics that better correlate with human judgments of quality~\cite{zhang2020bertscoreevaluatingtextgeneration,Fabbri2021SummEval}. However, there is little work in reproducing studies on evaluation metrics, despite our field's reliance on their outcomes. We include a number of traditional metrics in both our reproduction study and as baselines for our LLM-as-judge experiments. 

 Recent work has shown great success using LLMs-as-judges, showing high correlation with human judgments across a variety of tasks and domains~\cite{Zheng2023JudgingLLM}. Researchers have improved LLM-as-Judge performance through methods such as evaluation-specific finetuning~\cite{Kim2023Prometheus} or better prompting~\cite{liu-etal-2023-g}. Although LLMs-as-Judges perform best when provided with a gold reference~\cite{Deutsch2022LimitationsReferenceFree}, there is great interest in the reference-free setting of evaluation metrics as gold references are often unavailable or prohibitively expensive to collect~\cite{liu-etal-2023-g,vasilyev-etal-2020-fill}. We propose a new setting, reference-adjusted evaluation, in which relevant but non-exact references are available.